
% Auto-generated by figure/code/build_case_study_tex.py
% Required packages: graphicx, xcolor, array.
\definecolor{caseBlue}{HTML}{DDECF9}
\definecolor{caseGreen}{HTML}{198754}
\definecolor{caseRed}{HTML}{B42318}
\definecolor{caseGray}{HTML}{F3F4F6}
\newsavebox{\caseMeasureBox}
\newlength{\caseReasonWidth}
\newdimen\caseQwenH
\newdimen\caseSeedH
\newcommand{\caseTitle}[2]{%
  \noindent{\Large\bfseries #2 \texttt{#1}}%
  \par\vspace{0.45em}\hrule\vspace{0.55em}
}
\newcommand{\caseSubTitle}[1]{%
  {\bfseries #1}\par\vspace{0.2em}
}
\newcommand{\casePrompt}[1]{%
  {\footnotesize\emph{#1}}\par\vspace{0.35em}
}
\newcommand{\caseClip}[4]{%
  \begin{minipage}[t]{#4}
    \centering
    {\scriptsize\bfseries #1}\par\vspace{0.12em}
    \includegraphics[width=0.49\linewidth]{#2}\hfill
    \includegraphics[width=0.49\linewidth]{#3}\\[-0.1em]
    {\tiny first\hfill last}
  \end{minipage}
}
\newcommand{\caseReasonBody}[3]{%
  \raggedright
  \emergencystretch=1em
  {\scriptsize\bfseries #1}\quad {\scriptsize #2}\par
  {\scriptsize\emph{Reason:} #3}
}
\newcommand{\caseMeasureReason}[4]{%
  \begingroup
  \setlength{\fboxsep}{1.5pt}%
  \setlength{\caseReasonWidth}{0.31005\linewidth}%
  \sbox{\caseMeasureBox}{%
    \begin{minipage}{\dimexpr\caseReasonWidth-2\fboxsep-2\fboxrule\relax}
      \caseReasonBody{#2}{#3}{#4}
    \end{minipage}%
  }%
  \ifdim\dimexpr\ht\caseMeasureBox+\dp\caseMeasureBox\relax>#1
    \global#1=\dimexpr\ht\caseMeasureBox+\dp\caseMeasureBox\relax
  \fi
  \endgroup
}
\newcommand{\caseReason}[4]{%
  \vspace{0.28em}
  \begingroup
  \setlength{\fboxsep}{1.5pt}%
  \setlength{\caseReasonWidth}{0.975\linewidth}%
  \noindent\makebox[\linewidth][c]{%
  \fcolorbox{black!15}{caseGray}{%
    \begin{minipage}[t][#1][t]{\dimexpr\caseReasonWidth-2\fboxsep-2\fboxrule\relax}
      \caseReasonBody{#2}{#3}{#4}
      \vfill
    \end{minipage}
  }}%
  \endgroup\par
}

\begin{figure*}[t]
\centering
\caseTitle{DAY1\_A2\_ALICE\_17130000}{egolife}
\caseQwenH=0pt
\caseSeedH=0pt
\caseMeasureReason{\caseQwenH}{Qwen3.5-397B-A17B}{\textbf{Ans.} D (True middle); \textbf{GT} D (True middle); \textcolor{caseGreen}{Correct}}{The last frame of the BEGINNING segment shows a person inside the car handing a bottle/bag to the camera person outside. The first frame of Candidate D shows this exact same action, continuing the sequence. The last frame of Candidate D shows the camera facing a group of people standing outside (girl in denim jacket, guy in blue shirt), which matches perfectly with the first frame of the END segment.}
\caseMeasureReason{\caseQwenH}{Qwen3.5-397B-A17B}{\textbf{Ans.} A (Expanded); \textbf{GT} D (True middle); \textcolor{caseRed}{Wrong}}{Candidate A starts with the exact view of legs and shadows outside the car door that ends the BEGINNING segment. It then shows the girl inside putting on her hat and exiting the vehicle, ending with a view of the group standing outside which perfectly matches the start of the END segment.}
\caseMeasureReason{\caseQwenH}{Qwen3.5-397B-A17B}{\textbf{Ans.} D (True middle); \textbf{GT} D (True middle); \textcolor{caseGreen}{Correct}}{Candidate D starts with the person inside the car handing a water bottle (matching the end of the BEGINNING segment) and ends with the group standing outside (matching the start of the END segment), showing the continuous action of exiting the vehicle and gathering outside.}
\caseMeasureReason{\caseSeedH}{Seed1.8-T}{\textbf{Ans.} C (True middle); \textbf{GT} C (True middle); \textcolor{caseGreen}{Correct}}{Candidate C directly continues from the last frame of BEGINNING, showing the person in the white shirt being helped out of the car, and smoothly transitions to the outdoor group scene that matches the start of END, making the sequence logical and continuous.}
\caseMeasureReason{\caseSeedH}{Seed1.8-T}{\textbf{Ans.} A (True middle); \textbf{GT} A (True middle); \textcolor{caseGreen}{Correct}}{Candidate A smoothly continues from the last frame of BEGINNING, showing the person exiting the car, putting on a cap, and transitioning to the group standing outside which matches the first frame of END, maintaining a logical temporal and action flow.}
\caseMeasureReason{\caseSeedH}{Seed1.8-T}{\textbf{Ans.} A (Repeat); \textbf{GT} C (True middle); \textcolor{caseRed}{Wrong}}{Candidate A continues directly from the BEGINNING's final frame, showing the person in the car being helped out, exiting, and the door being closed, which smoothly transitions to the END's outdoor group interaction scene. The other candidates either reverse the sequence or start at the END's scene, breaking the temporal flow.}
\advance\caseQwenH by 0.25em
\advance\caseSeedH by 0.25em
\begin{center}
\caseClip{BEGINNING: first / last}{case_study_frames/DAY1_A2_ALICE_17130000_before_first.pdf}{case_study_frames/DAY1_A2_ALICE_17130000_before_last.pdf}{0.36\linewidth}
\hspace{1.0em}
\caseClip{END: first / last}{case_study_frames/DAY1_A2_ALICE_17130000_after_first.pdf}{case_study_frames/DAY1_A2_ALICE_17130000_after_last.pdf}{0.36\linewidth}
\end{center}
\vspace{0.4em}
\begin{minipage}[t]{0.318\linewidth}
\centering
\caseSubTitle{Question S: Semantic}
\casePrompt{Choose the middle clip that connects the shared BEGINNING and END.}
{\setlength{\tabcolsep}{0pt}\begin{tabular}{@{}c@{\hspace{0.02\linewidth}}c@{}}
\caseClip{True middle}{case_study_frames/DAY1_A2_ALICE_17130000_S_GT_first.pdf}{case_study_frames/DAY1_A2_ALICE_17130000_S_GT_last.pdf}{0.47\linewidth} & \caseClip{Rand1}{case_study_frames/DAY1_A2_ALICE_17130000_S_S_Rand1_first.pdf}{case_study_frames/DAY1_A2_ALICE_17130000_S_S_Rand1_last.pdf}{0.47\linewidth}\\[0.45em]
\caseClip{Rand2}{case_study_frames/DAY1_A2_ALICE_17130000_S_S_Rand2_first.pdf}{case_study_frames/DAY1_A2_ALICE_17130000_S_S_Rand2_last.pdf}{0.47\linewidth} & \caseClip{Rand3}{case_study_frames/DAY1_A2_ALICE_17130000_S_S_Rand3_first.pdf}{case_study_frames/DAY1_A2_ALICE_17130000_S_S_Rand3_last.pdf}{0.47\linewidth}
\end{tabular}}
\vspace{0.35em}
\caseReason{\caseQwenH}{Qwen3.5-397B-A17B}{\textbf{Ans.} D (True middle); \textbf{GT} D (True middle); \textcolor{caseGreen}{Correct}}{The last frame of the BEGINNING segment shows a person inside the car handing a bottle/bag to the camera person outside. The first frame of Candidate D shows this exact same action, continuing the sequence. The last frame of Candidate D shows the camera facing a group of people standing outside (girl in denim jacket, guy in blue shirt), which matches perfectly with the first frame of the END segment.}
\caseReason{\caseSeedH}{Seed1.8-T}{\textbf{Ans.} C (True middle); \textbf{GT} C (True middle); \textcolor{caseGreen}{Correct}}{Candidate C directly continues from the last frame of BEGINNING, showing the person in the white shirt being helped out of the car, and smoothly transitions to the outdoor group scene that matches the start of END, making the sequence logical and continuous.}
\end{minipage}
\hfill
\begin{minipage}[t]{0.318\linewidth}
\centering
\caseSubTitle{Question A: Alignment}
\casePrompt{Choose the middle clip that connects the shared BEGINNING and END.}
{\setlength{\tabcolsep}{0pt}\begin{tabular}{@{}c@{\hspace{0.02\linewidth}}c@{}}
\caseClip{True middle}{case_study_frames/DAY1_A2_ALICE_17130000_A_GT_first.pdf}{case_study_frames/DAY1_A2_ALICE_17130000_A_GT_last.pdf}{0.47\linewidth} & \caseClip{Advanced}{case_study_frames/DAY1_A2_ALICE_17130000_A_A_Early_first.pdf}{case_study_frames/DAY1_A2_ALICE_17130000_A_A_Early_last.pdf}{0.47\linewidth}\\[0.45em]
\caseClip{Deferred}{case_study_frames/DAY1_A2_ALICE_17130000_A_A_Late_first.pdf}{case_study_frames/DAY1_A2_ALICE_17130000_A_A_Late_last.pdf}{0.47\linewidth} & \caseClip{Expanded}{case_study_frames/DAY1_A2_ALICE_17130000_A_A_Wide_first.pdf}{case_study_frames/DAY1_A2_ALICE_17130000_A_A_Wide_last.pdf}{0.47\linewidth}
\end{tabular}}
\vspace{0.35em}
\caseReason{\caseQwenH}{Qwen3.5-397B-A17B}{\textbf{Ans.} A (Expanded); \textbf{GT} D (True middle); \textcolor{caseRed}{Wrong}}{Candidate A starts with the exact view of legs and shadows outside the car door that ends the BEGINNING segment. It then shows the girl inside putting on her hat and exiting the vehicle, ending with a view of the group standing outside which perfectly matches the start of the END segment.}
\caseReason{\caseSeedH}{Seed1.8-T}{\textbf{Ans.} A (True middle); \textbf{GT} A (True middle); \textcolor{caseGreen}{Correct}}{Candidate A smoothly continues from the last frame of BEGINNING, showing the person exiting the car, putting on a cap, and transitioning to the group standing outside which matches the first frame of END, maintaining a logical temporal and action flow.}
\end{minipage}
\hfill
\begin{minipage}[t]{0.318\linewidth}
\centering
\caseSubTitle{Question P: Progression}
\casePrompt{Choose the middle clip that connects the shared BEGINNING and END.}
{\setlength{\tabcolsep}{0pt}\begin{tabular}{@{}c@{\hspace{0.02\linewidth}}c@{}}
\caseClip{True middle}{case_study_frames/DAY1_A2_ALICE_17130000_C_GT_first.pdf}{case_study_frames/DAY1_A2_ALICE_17130000_C_GT_last.pdf}{0.47\linewidth} & \caseClip{Reverse}{case_study_frames/DAY1_A2_ALICE_17130000_C_C_Reverse_first.pdf}{case_study_frames/DAY1_A2_ALICE_17130000_C_C_Reverse_last.pdf}{0.47\linewidth}\\[0.45em]
\caseClip{Reorder}{case_study_frames/DAY1_A2_ALICE_17130000_C_C_Shuffle_first.pdf}{case_study_frames/DAY1_A2_ALICE_17130000_C_C_Shuffle_last.pdf}{0.47\linewidth} & \caseClip{Repeat}{case_study_frames/DAY1_A2_ALICE_17130000_C_C_Loop_first.pdf}{case_study_frames/DAY1_A2_ALICE_17130000_C_C_Loop_last.pdf}{0.47\linewidth}
\end{tabular}}
\vspace{0.35em}
\caseReason{\caseQwenH}{Qwen3.5-397B-A17B}{\textbf{Ans.} D (True middle); \textbf{GT} D (True middle); \textcolor{caseGreen}{Correct}}{Candidate D starts with the person inside the car handing a water bottle (matching the end of the BEGINNING segment) and ends with the group standing outside (matching the start of the END segment), showing the continuous action of exiting the vehicle and gathering outside.}
\caseReason{\caseSeedH}{Seed1.8-T}{\textbf{Ans.} A (Repeat); \textbf{GT} C (True middle); \textcolor{caseRed}{Wrong}}{Candidate A continues directly from the BEGINNING's final frame, showing the person in the car being helped out, exiting, and the door being closed, which smoothly transitions to the END's outdoor group interaction scene. The other candidates either reverse the sequence or start at the END's scene, breaking the temporal flow.}
\end{minipage}
\caption{Qualitative case study on egolife sample \texttt{DAY1\_A2\_ALICE\_17130000}. Each subtask shows the shared context, candidate middle clips, model answers, and model reasons.}
\label{fig:case-study-egolife-day1-a2-alice-17130000}
\end{figure*}

\begin{figure*}[t]
\centering
\caseTitle{HwAPGFLvoHg.60\_0}{lvd}
\caseQwenH=0pt
\caseSeedH=0pt
\caseMeasureReason{\caseQwenH}{Qwen3.5-397B-A17B}{\textbf{Ans.} D (Rand2); \textbf{GT} B (True middle); \textcolor{caseRed}{Wrong}}{Candidate D shows the knife being inserted into the pocket, which seamlessly continues the action from the BEGINNING segment where the knife is positioned at the pocket edge. It leads logically into the END segment where the knife is fully inserted and the person's hand is resting, before a pen is introduced.}
\caseMeasureReason{\caseQwenH}{Qwen3.5-397B-A17B}{\textbf{Ans.} A (True middle); \textbf{GT} A (True middle); \textcolor{caseGreen}{Correct}}{Candidate A shows the knife being fully inserted into the pocket and the hands smoothing the fabric, which seamlessly connects the beginning (inserting the knife) to the end (knife in pocket, hand resting).}
\caseMeasureReason{\caseQwenH}{Qwen3.5-397B-A17B}{\textbf{Ans.} B (True middle); \textbf{GT} B (True middle); \textcolor{caseGreen}{Correct}}{The beginning segment shows the knife being inserted into the pocket. Candidate B starts with the hands adjusting the knife in the pocket, which naturally follows the insertion. It ends with the hands smoothing the fabric and the right hand resting on the pants. This matches perfectly with the start of the end segment, where the right hand is resting on the pants before a pen is introduced.}
\caseMeasureReason{\caseSeedH}{Seed1.8-T}{\textbf{Ans.} A (True middle); \textbf{GT} A (True middle); \textcolor{caseGreen}{Correct}}{Candidate A smoothly transitions from the end of the BEGINNING sequence, where the hands are inserting the black tool into the pocket, to the start of the END sequence, where the tool is partially tucked into the pocket, matching the action flow perfectly.}
\caseMeasureReason{\caseSeedH}{Seed1.8-T}{\textbf{Ans.} C (True middle); \textbf{GT} C (True middle); \textcolor{caseGreen}{Correct}}{Candidate C shows a smooth transition from the BEGINNING's action of inserting the black tool into the pocket, adjusting it to the partially visible position that matches the first frame of the END, creating a continuous sequence of actions.}
\caseMeasureReason{\caseSeedH}{Seed1.8-T}{\textbf{Ans.} A (True middle); \textbf{GT} A (True middle); \textcolor{caseGreen}{Correct}}{Candidate A continues the action from the beginning, showing the hands pushing the black object fully into the pocket, which directly leads to the first frame of the END where the object is partially in the pocket with a hand resting on it. The sequence of actions is logically and visually consistent.}
\advance\caseQwenH by 0.25em
\advance\caseSeedH by 0.25em
\begin{center}
\caseClip{BEGINNING: first / last}{case_study_frames/HwAPGFLvoHg.60_0_before_first.pdf}{case_study_frames/HwAPGFLvoHg.60_0_before_last.pdf}{0.36\linewidth}
\hspace{1.0em}
\caseClip{END: first / last}{case_study_frames/HwAPGFLvoHg.60_0_after_first.pdf}{case_study_frames/HwAPGFLvoHg.60_0_after_last.pdf}{0.36\linewidth}
\end{center}
\vspace{0.4em}
\begin{minipage}[t]{0.318\linewidth}
\centering
\caseSubTitle{Question S: Semantic}
\casePrompt{Choose the middle clip that connects the shared BEGINNING and END.}
{\setlength{\tabcolsep}{0pt}\begin{tabular}{@{}c@{\hspace{0.02\linewidth}}c@{}}
\caseClip{True middle}{case_study_frames/HwAPGFLvoHg.60_0_S_GT_first.pdf}{case_study_frames/HwAPGFLvoHg.60_0_S_GT_last.pdf}{0.47\linewidth} & \caseClip{Rand1}{case_study_frames/HwAPGFLvoHg.60_0_S_S_Rand1_first.pdf}{case_study_frames/HwAPGFLvoHg.60_0_S_S_Rand1_last.pdf}{0.47\linewidth}\\[0.45em]
\caseClip{Rand2}{case_study_frames/HwAPGFLvoHg.60_0_S_S_Rand2_first.pdf}{case_study_frames/HwAPGFLvoHg.60_0_S_S_Rand2_last.pdf}{0.47\linewidth} & \caseClip{Rand3}{case_study_frames/HwAPGFLvoHg.60_0_S_S_Rand3_first.pdf}{case_study_frames/HwAPGFLvoHg.60_0_S_S_Rand3_last.pdf}{0.47\linewidth}
\end{tabular}}
\vspace{0.35em}
\caseReason{\caseQwenH}{Qwen3.5-397B-A17B}{\textbf{Ans.} D (Rand2); \textbf{GT} B (True middle); \textcolor{caseRed}{Wrong}}{Candidate D shows the knife being inserted into the pocket, which seamlessly continues the action from the BEGINNING segment where the knife is positioned at the pocket edge. It leads logically into the END segment where the knife is fully inserted and the person's hand is resting, before a pen is introduced.}
\caseReason{\caseSeedH}{Seed1.8-T}{\textbf{Ans.} A (True middle); \textbf{GT} A (True middle); \textcolor{caseGreen}{Correct}}{Candidate A smoothly transitions from the end of the BEGINNING sequence, where the hands are inserting the black tool into the pocket, to the start of the END sequence, where the tool is partially tucked into the pocket, matching the action flow perfectly.}
\end{minipage}
\hfill
\begin{minipage}[t]{0.318\linewidth}
\centering
\caseSubTitle{Question A: Alignment}
\casePrompt{Choose the middle clip that connects the shared BEGINNING and END.}
{\setlength{\tabcolsep}{0pt}\begin{tabular}{@{}c@{\hspace{0.02\linewidth}}c@{}}
\caseClip{True middle}{case_study_frames/HwAPGFLvoHg.60_0_A_GT_first.pdf}{case_study_frames/HwAPGFLvoHg.60_0_A_GT_last.pdf}{0.47\linewidth} & \caseClip{Advanced}{case_study_frames/HwAPGFLvoHg.60_0_A_A_Early_first.pdf}{case_study_frames/HwAPGFLvoHg.60_0_A_A_Early_last.pdf}{0.47\linewidth}\\[0.45em]
\caseClip{Deferred}{case_study_frames/HwAPGFLvoHg.60_0_A_A_Late_first.pdf}{case_study_frames/HwAPGFLvoHg.60_0_A_A_Late_last.pdf}{0.47\linewidth} & \caseClip{Expanded}{case_study_frames/HwAPGFLvoHg.60_0_A_A_Wide_first.pdf}{case_study_frames/HwAPGFLvoHg.60_0_A_A_Wide_last.pdf}{0.47\linewidth}
\end{tabular}}
\vspace{0.35em}
\caseReason{\caseQwenH}{Qwen3.5-397B-A17B}{\textbf{Ans.} A (True middle); \textbf{GT} A (True middle); \textcolor{caseGreen}{Correct}}{Candidate A shows the knife being fully inserted into the pocket and the hands smoothing the fabric, which seamlessly connects the beginning (inserting the knife) to the end (knife in pocket, hand resting).}
\caseReason{\caseSeedH}{Seed1.8-T}{\textbf{Ans.} C (True middle); \textbf{GT} C (True middle); \textcolor{caseGreen}{Correct}}{Candidate C shows a smooth transition from the BEGINNING's action of inserting the black tool into the pocket, adjusting it to the partially visible position that matches the first frame of the END, creating a continuous sequence of actions.}
\end{minipage}
\hfill
\begin{minipage}[t]{0.318\linewidth}
\centering
\caseSubTitle{Question P: Progression}
\casePrompt{Choose the middle clip that connects the shared BEGINNING and END.}
{\setlength{\tabcolsep}{0pt}\begin{tabular}{@{}c@{\hspace{0.02\linewidth}}c@{}}
\caseClip{True middle}{case_study_frames/HwAPGFLvoHg.60_0_C_GT_first.pdf}{case_study_frames/HwAPGFLvoHg.60_0_C_GT_last.pdf}{0.47\linewidth} & \caseClip{Reverse}{case_study_frames/HwAPGFLvoHg.60_0_C_C_Reverse_first.pdf}{case_study_frames/HwAPGFLvoHg.60_0_C_C_Reverse_last.pdf}{0.47\linewidth}\\[0.45em]
\caseClip{Reorder}{case_study_frames/HwAPGFLvoHg.60_0_C_C_Shuffle_first.pdf}{case_study_frames/HwAPGFLvoHg.60_0_C_C_Shuffle_last.pdf}{0.47\linewidth} & \caseClip{Repeat}{case_study_frames/HwAPGFLvoHg.60_0_C_C_Loop_first.pdf}{case_study_frames/HwAPGFLvoHg.60_0_C_C_Loop_last.pdf}{0.47\linewidth}
\end{tabular}}
\vspace{0.35em}
\caseReason{\caseQwenH}{Qwen3.5-397B-A17B}{\textbf{Ans.} B (True middle); \textbf{GT} B (True middle); \textcolor{caseGreen}{Correct}}{The beginning segment shows the knife being inserted into the pocket. Candidate B starts with the hands adjusting the knife in the pocket, which naturally follows the insertion. It ends with the hands smoothing the fabric and the right hand resting on the pants. This matches perfectly with the start of the end segment, where the right hand is resting on the pants before a pen is introduced.}
\caseReason{\caseSeedH}{Seed1.8-T}{\textbf{Ans.} A (True middle); \textbf{GT} A (True middle); \textcolor{caseGreen}{Correct}}{Candidate A continues the action from the beginning, showing the hands pushing the black object fully into the pocket, which directly leads to the first frame of the END where the object is partially in the pocket with a hand resting on it. The sequence of actions is logically and visually consistent.}
\end{minipage}
\caption{Qualitative case study on lvd sample \texttt{HwAPGFLvoHg.60\_0}. Each subtask shows the shared context, candidate middle clips, model answers, and model reasons.}
\label{fig:case-study-lvd-hwapgflvohg.60-0}
\end{figure*}
